% Back matter. The statements are factual records of how the work was produced.

\section*{Declarations}

\bmhead{Use of generative AI}
This work was carried out by a generative AI system, listed as first author, under human
direction, and we describe the division of work in full. The software
described in this paper (the engine-independent library, the three engine backends, the
scenarios, the baselines and the tools), the design of the benchmark, the analysis of its
results and the draft of this manuscript were produced with Claude (Anthropic), used through
its coding agent Claude Code, under the direction of the second author. The work, from 10 to 12 November 2026, was carried out mostly by the model
Claude Fable~5.2, which also directed further instances of itself, of Claude Opus~5.5 (an
independent audit of the code and referee reports) and of Claude Sonnet~5.5 (auxiliary tasks:
scripts, packaging, the viewer and the site). The
starting point was the second author's 2018 thesis \cite{Vilela2018} and its code. The second author set the
questions, decided the scope, approved the changes to the contact laws and the rules of the
benchmark. Under that direction the system reviewed the literature and
proposed the framing of the paper as a benchmark; designed and implemented the laws' evaluation,
including the derivation and the numerical check of the formulas of Section~\ref{sec:friction};
audited the first version of the benchmark and redesigned it, including the classification of
the references and the second set of scenarios; configured the native models from the engines'
documentation; ran the battery; wrote the scripts that compute the statistics and generate
every number, table and figure of the paper from the result files; drafted the text; and, in a
further session, wrote a referee's report on the manuscript and revised the analysis and the
text in response to it (12 November 2026). The
schematic figures are diagrams drawn by the system from the geometry of the paper and the plots are generated
from the data; there are no generated images. Every reference in the bibliography was checked
against the publisher's or registry's record, and the repository records which sources were
read in full. The first author is an AI system and cannot be held accountable; the second author
directed the work and is the person to whom questions about it should be addressed.

\bmhead{Funding}
No funding was received for this work.

\bmhead{Competing interests}
The second author has no competing interests to declare; the first author is a commercial AI
system developed by Anthropic.

\bmhead{Data and code availability}
The library, the three engine backends, the scenarios with their references, the mesh generator,
the scripts that run the battery and compute its statistics, and the sources of this paper with
the scripts that generate its numbers, tables and figures and check its derivations and
references are open source under the BSD 3-Clause licence at
\code{https://github.com/dvilelaf/volumetric-contact}. The
build runs in a container with pinned versions of Drake, Project Chrono and MuJoCo. The results
of the battery reported here (\code{battery.csv}, with checksum \result{battery/sha256};
\code{terms.csv}; \code{stats.json}) are part of the repository, one command regenerates from
them every number, table and figure of the paper, and the trajectories of the runs can be
replayed in a viewer published with it. The main matrix was run on \result{battery/date} with
the code of commit \result{battery/commit}; the friction quadratures and \law{pcm} were added
at commit \result{battery/commit-baselines}, after \result{battery/spot-runs} runs spread over
the scenarios reproduced the main matrix (largest relative difference
\result{battery/spot-diff}); the timings are those of commit \result{battery/timing-commit} on
\result{battery/timing-date}. The derivations of the predictions quoted in the text and the
single diagnostic runs behind them are recorded in the history of the repository.

\bmhead{Author contributions}
Claude Fable~5.2: methodology, software, validation, formal analysis, investigation, data
curation, visualization, writing (original draft, review and editing). David Vilela:
conceptualization, supervision, project administration, and the 2018 thesis from which the work
derives \cite{Vilela2018}.

\bmhead{Ethics approval, consent to participate, consent for publication}
Not applicable.
